\chapter{A két adatkészlet empirikus összehasonlítása: flexset versus soloq}

\section{Az összehasonlítás módszertani szerepe}

A projekt egyik legfontosabb kutatási-módszertani döntése az volt, hogy nem egyetlen adatkészlettel dolgozik, hanem egymással összehasonlítható, eltérő queue-szemantikájú korpuszokkal. Ez nem pusztán mennyiségi bővítés, hanem kvalitatív különbségtétel: a \textit{flexset} (EUNE Master$+$ Flex Queue) és a \textit{soloq} (EUNE Master$+$ Solo/Duo Queue) nemcsak méretükben, hanem a bennük keletkező kapcsolatok jellegében is eltérnek.

A kutatási kérdés így fogalmazható meg: ha ugyanazon analitikai pipeline-t lefuttatjuk mindkét dataseten, mikor kapunk eltérő eredményt, és ez az eltérés értelmezhető-e a két queue eltérő mechanikájából? Az egyetlen datasetre épülő eredmény interpretálása mindig kétséges: az eredmény lehet a módszer hatása, lehet a dataset sajátossága, vagy lehet valódi gráfstruktúra. Két, eltérő generáló mechanizmussal bíró dataseten kapott konvergens vagy divergens eredmény sokkal erősebb támpontot ad.

\section{A két dataset konstruktív különbségei}

\subsection{Queue-szemantika és kapcsolatgeneráló mechanizmus}

A Flex Queue (440-es queue ID) csapatszedést tesz lehetővé: játékosok 2--5 főig premade csapatot alkothatnak, és egységként lépnek be a meccsbe. Ez azt jelenti, hogy a szövetségesi pozícióban repeatedly megjelenő játékospárok egy része \emph{szándékolt ismétlődés}: a premade csapat azonos tagjai meccsről meccsre együtt indulnak. Az ally/enemy arány (73.17\% ally a \textit{flexset}-en) részben ennek a mechanizmusnak a következménye.

A Solo/Duo Queue (420-as queue ID) ezzel szemben legfeljebb két fős premade-et engedélyez. A matchmaking itt jóval véletlenszerűbb: egy adott játékos öttagú csapatának többi tagja döntően nem szándékolt, hanem algoritmikus összerendelés. Az ally/enemy arány (65.24\% ally a \textit{soloq}-on) alacsonyabb, ami összhangban van azzal, hogy a véletlenszerűbb meccsek kevesebb domináns szövetségesi kapcsolatot hoznak létre.

\subsection{Adatskála és meccssűrűség}

\begin{table}[H]
\centering
\caption{A két dataset alapvető méretei}
\begin{tabularx}{\textwidth}{>{\raggedright\arraybackslash}p{5cm} rr}
\toprule
\textbf{Mutató} & \textbf{flexset} & \textbf{soloq} \\
\midrule
Meccs-JSON állományok & 4981 & 2001 \\
Nyers játékosok & 15\,147 & 2822 \\
Finomított játékosok & 15\,421 & 6768 \\
Klaszterek (DB) & 599 & 140 \\
Klasztertagok (DB) & 5741 & 1538 \\
Graph V2 látható csúcsok & 5741 & 1538 \\
Graph V2 látható élek & 18\,548 & 4111 \\
Átlagos élsúly & 3.405 & 2.364 \\
Graph V2 ally/enemy arány & 73.17\% / 26.83\% & 65.24\% / 34.76\% \\
Graph V2 klaszterek & 1531 & 870 \\
Átlagos klaszterméret & 3.750 & 1.768 \\
Legnagyobb klaszterméret & 20 & 20 \\
\bottomrule
\end{tabularx}
\end{table}

Az átlagos élsúly különbsége (3.405 vs. 2.364) különösen figyelemreméltó. Ez azt jelenti, hogy a \textit{flexset}-en a látható kapcsolatok átlagosan sűrűbb, ismétlődőbb előfordulásból épülnek fel. Ahol a \textit{soloq}-on egy él két meccs ismétlődéséből áll össze, ott a \textit{flexset}-en tipikusan több mint három meccs van mögötte. Ez közvetlen mérőszáma annak, hogy a Flex Queue-ban az azonos játékospárok egymással sűrűbben, tartósabban találkoznak.

\section{A két dataset vizuális klaszterstruktúrájának különbsége}

Az átlagos klaszterméret különbsége (3.750 vs. 1.768) szintén tartalmas. A \textit{flexset}-en a csoportok átlagosan majdnem négytagúak, míg a \textit{soloq}-on szinte mindenütt kis, kétfős vagy egyedülálló klaszterek dominálnak. Az 1.768-as átlag azt jelenti, hogy a \textit{soloq}-ban a Graph V2 által látható csomópontok közel fele egyedül álló csúcs, azaz nincs mellette kellő súlyú szövetségesi partner.

Ez a különbség nem adatgyűjtési hiba, hanem a queue-mechanika közvetlen következménye. A premade Flex Queue-ban természetesen keletkeznek 3--5 fős, ismétlődő klasztermagok, amelyek a Graph V2-ben összefüggő csoportként jelennek meg. A véletlenszerűbb \textit{soloq}-ban ez a csoportosulás nem jön létre organikusan.

A \textit{flexset}-en 815 bridgejelölt-klaszter van 1531 összesből (53.2\%), a \textit{soloq}-on 605 jelölt 870-ből (69.5\%). Ez meglepő arány: a \textit{soloq} klaszterei relatíve több hidat képeznek. Ennek magyarázata valószínűleg az, hogy a \textit{soloq}-on a kis klaszterek szétszórtabb, kevésbé összefüggő grafikonban helyezkednek el, így relatív módon több közöttük a vizuálisan kitüntetett pozícióban lévő.

\section{A metrikafedettség összehasonlítása}

Mindkét dataseten a Graph V2 artifact teljesen fedett \textit{opscore} és \textit{feedscore} értékekkel. A \textit{flexset}-en az \textit{opscore}-fedettség 5741/5741, a \textit{soloq}-on 1538/1538; ugyanez igaz a \textit{feedscore}-ra. A teljes fedettség azt jelenti, hogy a Graph V2-ben megjelenő minden látható csomóponthoz rendelkezésre áll érvényes metrika, és nincs olyan csomópont, amely hiányzó score miatt kiesne az elemzésből. Az \textit{opscore} átlagok közel esnek egymáshoz (\textit{flexset}: 6.1015, \textit{soloq}: 6.1701), ami megerősíti, hogy az összehasonlítás az összehasonlíthatóság feltételei között zajlik.

\section{A teljes összehasonlítás tanulságai}

\begin{table}[H]
\centering
\caption{A két dataset összehasonlításának összefoglalója}
\begin{tabularx}{\textwidth}{>{\raggedright\arraybackslash}p{5cm} X X}
\toprule
\textbf{Dimenzió} & \textbf{flexset} & \textbf{soloq} \\
\midrule
Queue-mechanika & Flex, premade-fókusz & Solo/Duo, egyénközpontú \\
Átlagos élsúly & 3.405 (sűrűbb) & 2.364 (ritkább) \\
Ally arány & 73.17\% & 65.24\% \\
Átlagos klaszterméret & 3.750 & 1.768 \\
Triádok/csomópont & 4.95 & 1.03 \\
Core-periphery olvasat & Hídgazdag asszociatív mag és sok periférikus ally-sziget & Kisebb, véletlenszerűbb lokális csoportok \\
Social-path opscore assz. & 0.2346 & 0.1750 \\
Battle-path opscore assz. & 0.1461 & 0.1625 \\
Feedscore battle-path assz. & $-0.0108$ & $-0.0366$ \\
Kutatási szerep & Asszociatív gráf és teljesítménykapcsolat & Teljesítmény-alapvonal, kontroll \\
\bottomrule
\end{tabularx}
\end{table}

Az összehasonlítás általános tanulsága: a \textit{flexset} az asszociatív játékosgráf és a visszatérő szövetségesi kapcsolatok természetesebb terepe, a \textit{soloq} pedig a teljesítmény-alapvonal és a kontroll-dataset szerepét tölti be. Ahol a két dataset hasonló eredményt ad (mindkettőn pozitív social-path \textit{opscore} asszortativitás), ott az eredmény általánosabb jellegű. Ahol eltérnek (klaszterméret, élsűrűség, bridge-orbit szerkezet), ott az eltérés a queue-mechanikából fakad, és önmagában is tudományosan értelmezhető.

Figyelemreméltó részlet: a \textit{battle-path opscore} értéke a \textit{soloq}-on magasabb (0.1625), mint a \textit{flexset}-en (0.1461). Ennek magyarázata a \textit{soloq} szűk MMR-sávján alapuló matchmaking: az ismétlődően egymással találkozó játékosok \textit{opscore}-ban közel esnek egymáshoz, de ez a struktúra gyengébb és esetlegesebb, mint a \textit{flexset} szándékolt premade-kapcsolatai. A \textit{social-path} módban a \textit{flexset} erősebb (0.2346 vs. 0.1750), ami a premade-szövetségesi mechanizmust tükrözi.
